\documentclass[11pt]{article}
\usepackage[margin=1in]{geometry}
\usepackage{amsmath,amssymb,amsthm,booktabs,hyperref,microtype}
\hypersetup{colorlinks=true,linkcolor=blue,pdftitle={Robust Multichannel Perron Envelopes},pdfauthor={Yang--Mills representational dynamics research programme}}
\newtheorem{theorem}{Theorem}
\newtheorem{corollary}{Corollary}
\newtheorem{warning}{Warning}
\title{Robust Multichannel Perron Envelopes\\Target 56 of the H1 Downstream Programme}
\author{Yang--Mills representational dynamics research programme}
\date{21 July 2026}
\begin{document}
\maketitle
\begin{abstract}
We replace a single calibrated defect kernel by a robust entrywise envelope.
The Perron vector of the envelope controls every dominated physical candidate.
For the four-channel H1 star we calculate the exact additive, multiplicative,
and symmetric persistence/tunneling slack before the degree-18 cluster
threshold.  The optimization problem is closed; constructing the complete
nonlinear channel envelopes remains open.
\end{abstract}

\section{Envelope theorem}
Let $Q,B$ be finite nonnegative matrices with $Q\leq B$ entrywise.  Suppose
$B$ is irreducible and $BV_B=\rho(B)V_B$ with $V_B>0$.
\begin{theorem}[Robust Perron domination]
\[
 QV_B\leq \rho(B)V_B.
\]
Consequently every kernel in the order interval $[0,B]$ has weighted ceiling
at most $\rho(B)$ in the common weight $V_B$.
\end{theorem}
\begin{proof}
Positivity gives $QV_B\leq BV_B$ componentwise.  The conclusion follows from
the Perron equation.
\end{proof}

This theorem is deliberately one-sided.  A numerically plausible row not
proved to lie below $B$ receives no certificate.

\section{Four-channel uncertainty budget}
Let $k=4$ and use the root/charged matrix
\[
 B_{0a}=\frac{q^2}{k},\qquad B_{a0}=1-s,\qquad
 B_{aa}=p,\qquad B_{ab}=\frac{c}{k-1}\quad(a\neq b),
\]
where $s=p+c$.  On the symmetric charged subspace the Perron root is
\begin{equation}\label{eq:rho}
 \rho(s)=\frac{s+\sqrt{s^2+4q^2(1-s)}}2.
\end{equation}
For $s=0$, $\rho=q=0.002430277501841998$.  With
$q_{\rm crit}=0.0027927080803578997$, the exact budgets are
\begin{align*}
 q_{\rm crit}-q&=0.0003624305785159018,\\
 \frac{q_{\rm crit}}q-1&=0.1491313556748981,\\
 s_{\max}&=\frac{q_{\rm crit}^2-q^2}{q_{\rm crit}-q^2}
 =0.0006792623948561134.
\end{align*}
The last number is the total symmetric charged-to-charged persistence or
tunneling mass compatible with strict subcriticality.  At half this budget the
certified radius is approximately $0.00260560708<q_{\rm crit}$.

\begin{corollary}[Additive weighted errors]
If $BV\leq qV$ and $EV\leq\varepsilon V$, then
$(B+E)V\leq(q+\varepsilon)V$.  Hence every normalization, chart-complement,
and cross-channel error consumes the same additive weighted slack.
\end{corollary}

\section{Physical channel ledger}
The required envelope must include continuous incoherence, center frustration,
cut/orbit charts, compositionally complete integer topology, persistence, and
tunneling.  Entrywise domination is sufficient, but a directly weighted bound
may be sharper.  Either route must use normalized conditional rows uniformly
over positive-mass histories.

\begin{warning}
The finite envelope is not a physical Yang--Mills kernel.  Omitting a channel,
using an unnormalized numerator, or accepting equality at $q_{\rm crit}$
invalidates the cluster conclusion.
\end{warning}

\section{Audit}
The robust theorem and budgets are C2.  Eleven hostile tests verify spectral
formulas, threshold strictness, domination, cross-channel equivalence, invalid
weights, and superprobability rejection.
\end{document}
